\section{Experimental setup}
\label{sec:setup}

All numbers in this and the following sections are produced by the scripts in the \texttt{experiments/} directory of the public repository. Each script writes one JSON record per run, and the tables are generated from these records.

\subsection{Datasets}
\label{subsec:datasets}

\paragraph{Elliptic Bitcoin dataset} The Elliptic dataset \citep{weber2019anti} contains 203,769 Bitcoin transactions (nodes) and 234,355 directed payment flows (edges) in 49 time steps of about two weeks each. 4,545 transactions (2.2\%) are labelled illicit, 42,019 (20.6\%) licit and the remaining 157,205 (77.1\%) are unlabelled. Each transaction has 166 anonymised features. The first is the time step. The next 93 are \emph{local} features of the transaction (the dataset documentation mentions, among others, the numbers of inputs and outputs, fees and output volume). The last 72 are \emph{aggregated} features, statistics over the transaction's one-hop neighbourhood. The time step is not used as a feature, so every model receives $d=165$ features. No edge crosses a time step. We use a temporal split: steps 1--30 for training (26,905 labelled transactions, 2,954 illicit), steps 31--34 for validation (2,989; 508 illicit) and steps 35--49 for testing (16,670; 1,083 illicit). \ref{app:elliptic} (Table~\ref{tab:elliptic_steps}) lists the number of licit, illicit and unlabelled transactions in every time step.

\paragraph{T-Finance} T-Finance \citep{tang2022rethinking} is a transaction network of a financial company with 39,357 accounts, 21.2 million undirected edges (transactions between accounts) and 10 account features; 4.58\% of the accounts are labelled fraudulent by human experts. We use the version and the first five fully supervised splits (40\%/20\%/40\% train/validation/test) of GADBench \citep{tang2023gadbench}. T-Finance has no timestamps, so it tests whether SEAL and the explanation objective carry over to a second real-world financial network with different features, entities and label source. HRAPE is switched off and temporal baselines are not applicable. To keep full-batch training tractable in memory, graph models see at most 10 uniformly sampled neighbours per account, which reduces the graph to about 0.4 million undirected edges.

\paragraph{Synthetic temporal laundering benchmark} The synthetic benchmark contains a signal that is known by construction and that only a model using timing information can exploit. The generator (\texttt{rtxgnn/synthetic.py}) works as follows.
(1)~A Barab\'asi--Albert graph on $n=5{,}000$ accounts ($m=2$) defines which pairs transact. Each pair gets a random direction and a Poisson number of transactions (mean 3), with timestamps uniform on $[0,180]$ days and log-normal amounts ($\mu=5,\sigma=1$).
(2)~Fifty laundering rings are laid on disjoint accounts. Each is a directed cycle of length $L\sim\mathcal{U}\{4,\dots,8\}$ executed within a \emph{burst} of 2 days at a random start time, with a 1--3\% fee per hop. Ring members keep their ordinary transactions.
(3)~Fifty \emph{decoy} cycles with the same length and amount distribution are laid on other accounts, but their hops are spread over the whole horizon. The cycle structure alone therefore does not identify laundering; the timing of consecutive hops does.
(4)~Every account has 16 features: 7 behavioural statistics computed from all its transactions (in/out counts, mean and standard deviation of in/out log-amounts, number of counterparties; standardised), a constant, and 8 i.i.d.\ $\mathcal{N}(0,1)$ noise features.
Ring members form the positive class (6.1\% of accounts), and the ground-truth explanation of a ring member consists of its two ring transactions. All accounts are scored at $t=180$ days with a stratified random 60/20/20 split.
Besides this \emph{default} setting we generate four variants: \emph{no\_decoy} (no decoy cycles, so structure alone suffices), \emph{wide\_burst} (rings completed within 10 instead of 2 days), \emph{label\_noise} (5\% of labels flipped) and \emph{shift} (+0.5 added to the noise features of ring members, a weak feature signal). Every setting is generated with seeds 0--4, and seed $s$ is also the model seed. By default the features carry no class signal, so the task cannot be solved from the features alone.


\subsection{Preprocessing}

For all neural models, each feature is mapped to a standard normal distribution by a quantile transform fitted on the transactions of the training steps only (Elliptic: steps 1--30, labelled and unlabelled; T-Finance: training accounts). No labels are used in this fit. Elliptic features are heavy-tailed, and the transform improves every neural model substantially (Table~\ref{tab:preprocessing} in \ref{app:preprocessing} reports results with the raw features). The transform was adopted during development, when test results were also computed. To rule out a choice that fits only the test period, we checked it on the validation steps, which are used for model selection: the transform raises the validation AP of \preNvalBetter{} of the \preNmodels{} neural models and leaves that of the sixth (SEFraud) practically unchanged. Tree ensembles and logistic regression use the raw features. The transform does not help trees, which are invariant to monotone transformations, and on Elliptic its quantisation slightly reduced the random forest's score.

\subsection{Baselines}
\label{subsec:baselines}

We compare against 16 baselines in four groups. \ref{app:baselines} describes how each one is adapted.
\begin{itemize}
\item \textbf{Feature-only models:} logistic regression (LR, class-balanced), random forest (RF; 100 trees, 50 features per split, the configuration of \citealp{weber2019anti}) and XGBoost \citep{chen2016xgboost} (300 trees, depth 6, class-balanced), and a two-layer MLP.
\item \textbf{General GNNs:} GCN \citep{kipf2017gcn}, GraphSAGE \citep{hamilton2017graphsage} and GAT \citep{velickovic2018gat}.
\item \textbf{Temporal GNNs:} EvolveGCN-O \citep{pareja2020evolvegcn}, TGAT \citep{xu2020tgat}, TGN \citep{rossi2020tgn} and APAN \citep{wang2021apan}.
\item \textbf{Fraud-specific GNNs:} CARE-GNN \citep{dou2020caregnn}, PC-GNN \citep{liu2021pcgnn}, GAS \citep{li2019gas}, FRAUDRE \citep{zhang2021fraudre}, and the self-explaining SEFraud \citep{li2024sefraud}.
\end{itemize}
All graph baselines were re-implemented in the same code base, because the original implementations target multi-relation review graphs (CARE-GNN, PC-GNN, FRAUDRE), heterogeneous graphs (GAS, SEFraud) or link prediction on continuous-time graphs (TGAT, TGN, APAN). The re-implementations follow the published architectures, adapted to the single-relation Elliptic graph. We cannot guarantee that they reach the performance of the authors' own code in the original settings.

DySAT is not included because its temporal attention is vacuous on Elliptic (each node appears in a single snapshot; Section~\ref{sec:related}). H2-FDetector and DGA-GNN, which were compared in the first version of this study under a different protocol, are not included either. H2-FDetector separates homophilic from heterophilic connections within each relation of a multi-relation graph, and DGA-GNN groups non-additive raw attributes by decision-tree binning. On Elliptic's single-relation graph with engineered features, their central mechanisms are represented by the label-aware neighbour selection of CARE-GNN and PC-GNN, the node--neighbour differences of FRAUDRE, and the tree ensembles, which are all included.

\subsection{Protocol and statistical analysis}
\label{subsec:protocol}

\paragraph{Training} All neural models use the same protocol and budget: hidden size 64, two message-passing layers where applicable, dropout 0.2, the class-balanced loss of Eq.~\eqref{eq:pred_loss} (plus the method's own auxiliary losses), Adam with learning rate $5\cdot10^{-3}$ and weight decay $10^{-5}$, gradient clipping at norm 1, at most 300 full-batch epochs, and early stopping on validation average precision with a patience of 40 epochs. No model, including RTXGNN, was tuned on the test data, and no baseline received a per-model hyper-parameter search. This is a limitation that affects all methods equally.

\paragraph{Decision threshold} Every model outputs a score. The decision threshold is the one that maximises F1 on the validation set, and it is then applied unchanged to the test set. We also report threshold-free metrics: the area under the ROC curve (AUC) and average precision (AP, the area under the precision--recall curve). AP is the more informative of the two under class imbalance.

\paragraph{Metrics} We report F1, precision and recall of the illicit class, together with AUC and AP. The overall test F1 is computed \emph{once over all 16,670 labelled test transactions pooled together} (micro-averaging over time steps), as in \citet{weber2019anti}. It is not the mean of the per-step F1 scores. Section~\ref{subsec:temporal} explains why the two differ.

\paragraph{Seeds and tests} Each configuration on Elliptic is run with 10 random seeds (0--9), which change the parameter initialisation, dropout masks and, for the random forest and XGBoost, the bootstrap and column sampling. We report mean $\pm$ standard deviation over the seeds. Logistic regression is deterministic. We compare RTXGNN with every baseline with a two-sided Wilcoxon signed-rank test on the per-seed F1 and AP pairs \citep{demsar2006statistical}. Seeds are paired because both models use the same seed. We apply Holm's correction \citep{holm1979simple} over the 16 comparisons of each metric. Ablations use 5 seeds (3 for the secondary variants), T-Finance uses GADBench splits 0--4, and the synthetic benchmark uses 5 generated graphs per setting. With 5 or 3 pairs, a two-sided Wilcoxon test cannot reach $p<0.05$, so the ablations report paired $t$-tests. These $p$-values are not corrected for multiple comparisons and should be read as indicative.

\paragraph{Hardware} Experiments were distributed over two kinds of machines: a cloud virtual machine with 4 vCPUs of an Intel Xeon processor (\cpuName) and 15\,GB of memory (CPU only), and Google Colab virtual machines with an NVIDIA T4 GPU. In the main comparison, seeds 0--2 and 5--9 ran on the CPU machine (except RTXGNN with seed 2) and seeds 3--4 on the GPU. The primary ablation study, the synthetic benchmark and the explanation study ran on the CPU machine. The EvolveGCN runs, the label-efficiency, retraining and T-Finance experiments, and the secondary ablations ran on the GPU. Each secondary ablation is compared with a reference model trained on the same hardware. Because the studies use different seed sets and hardware, the F1 of the full RTXGNN differs between them: \rtxFone{} in the main comparison (10 seeds), \ablfullFone{} in the primary ablation (5 seeds, CPU), \ablGfullFone{} in the secondary ablation (3 seeds, GPU) and \lertxgnnOneZeroZero{} with all labels in the label-efficiency study (5 seeds, GPU). Each study is compared only with its own reference. Both machines used PyTorch~2 and PyTorch Geometric~2 without mixed precision, and the algorithm is identical on both. All latency measurements (Section~\ref{subsec:latency}) were taken on the CPU machine.
